\chapter{Instalación}

Este capítulo deja el entorno listo y comprueba que funciona. Cómo arrancar el orquestador está en el capítulo siguiente.

\section{Requisitos}

\begin{itemize}
  \item \texttt{uv}, que gestiona Python y las dependencias (D-12). Si falta Python 3.14, se instala con \texttt{uv python install 3.14} (D-13).
  \item Docker con Compose, para PostgreSQL (D-69).
  \item \texttt{latexmk} y LuaLaTeX, solo para compilar la documentación.
\end{itemize}

\section{Pasos}

Instala las dependencias desde \texttt{uv.lock}:

\begin{verbatim}
uv sync
\end{verbatim}

Levanta PostgreSQL y espera a que esté listo:

\begin{verbatim}
docker compose up -d --wait
\end{verbatim}

Comprueba la instalación:

\begin{verbatim}
uv run pytest
\end{verbatim}

Todos los tests deben pasar. Los que usan la base de datos fallan con \texttt{Connection refused} si PostgreSQL no está levantado.

\section{Base de datos}

El contenedor usa el puerto 55432 y solo acepta conexiones desde el propio equipo. Por eso no choca con un PostgreSQL local en el puerto 5432. Sus credenciales son solo para desarrollo.

Los tests usan la variable de entorno \texttt{DATABASE\_URL} si está definida y, si no, la conexión al contenedor.

Para parar el contenedor sin perder los datos:

\begin{verbatim}
docker compose stop
\end{verbatim}
